\documentclass{beamer}
\title
{Resolución numérica del problema no lineal de mínimos cuadrados.
Aplicaciones a las estimación de parámetros de modelos matemáticos.}

\author
{Dídac Blanco Morros}


\date
{14 de febrero de 2023}

\titlegraphic{\vspace{0.7cm}\hspace{1.5cm}\includegraphics[height=1cm]{logomath.pdf}}
%\logo{\includegraphics[height=1cm]{overleaf-logo}}

\AtBeginSection[]
{
  \begin{frame}
    \frametitle{Table of Contents}
    \tableofcontents[currentsection]
  \end{frame}
}

\begin{document}
\frame{\titlepage}

\begin{frame}
\frametitle{Table of Contents} % Table of contents title
\tableofcontents
\end{frame}

\section{Fundamentos de la optimización sin restricciones}
\begin{frame}
    \frametitle{Conceptos previos}
    \begin{itemize}
        \item Función objetivo
        \item Función de error
        \item Función de costo
        \item Función de pérdida
        \item Función de riesgo
    \end{itemize}
\end{frame}

\section{Mínimos Cuadrados}


\section{El método de Levenberg-Marquardt}


\section{Implementación en Matlab\textsuperscript{®}}
\end{document}